\begin{afterword}
\summaryhead

The post suggests that people confuse belief, truth and reality because they think of beliefs as all-or-nothing. An imagined postmodernist says that ``The Sun goes around the Earth'' is true for a hunter-gatherer, while the reverse is true for an astronomer; the author replies that societies differ in beliefs, not truths. The opponent then says that saying you believe snow is white and saying it is true that snow is white express the same opinion, and quotes Wittgenstein. The author answers that the two-valued contrast of belief and disbelief looks too much like that of true and false.

With probabilities the types separate. A belief is a probability between 0 and 1. Its accuracy is scored against reality by the logarithm of the probability it gave to what happened: a 70\% belief that snow is white scores $-0.51$ bits if snow is white and $-1.73$ if not, and $-0.88$ in expectation. A belief about one's own belief is usually near-certain. So it is a type error to call a probability assignment ``true''. Keeping these levels distinct, the author hopes, makes the Mind Projection Fallacy less likely: a coin is not itself 50\% uncertain.

\responsehead

The quantitative part of the post is correct. The log score is a standard proper scoring rule, the three numbers are correct, and the actual score does always differ from the expected one.

The opponent at the start is ``the archetypal postmodernist'', and no postmodernist is named or quoted. The post answers the version it wrote. Its first reply, that societies differ in beliefs rather than truths, is enough against that version; the reader does not learn whether any actual defender of relativism says what the imagined postmodernist says.

The Wittgenstein line is quoted accurately. In the same section of the \textsc{Philosophical Investigations}, Wittgenstein says that ``I believe that this is the case'' is used like the assertion ``This is the case'', and yet that the hypothesis that I believe it is not used like the hypothesis that it is so. That is close to both sides of the exchange: the opponent's claim and the post's reply.

The central claim is causal: all-or-nothing thinking is ``a primary cause'' of the confusion, and thinking in probabilities makes it less likely. The post offers both with hedges (``I suggest'', ``I hope'', ``perhaps''), and gives no evidence for either. It shows that the distinctions can be drawn in a probabilistic framework, which is true. It does not show that people who use the framework confuse belief and truth less, and the hedges say as much.

In short: a correct and clear account of scoring beliefs against reality, framed by a postmodernist opponent whom the post does not source, with its claim about the cause and cure of the confusion offered as a suggestion and a hope, without evidence.
\end{afterword}
